\section{The ETR bridge}
\label{sec:etr-reduction}

To obtain \(\ExistsR\)-hardness, a reduction must begin with an explicitly
encoded ETR sentence
\begin{equation}
  \varphi=\exists x_1,\ldots,x_m:\Psi(x_1,\ldots,x_m)
\end{equation}
and construct in polynomial time an instance \(I_\varphi\) of the exact
language in \cref{sec:problem} such that
\begin{equation}
  \varphi\text{ is true}
  \quad\Longleftrightarrow\quad
  \Feas(I_\varphi).
  \label{eq:etr-reduction-spec}
\end{equation}

Related neural-network training languages are known to be
\(\ExistsR\)-complete~\cite{AbrahamsenKleistMiltzow2021}.  A citation to that
result does not by itself prove \cref{eq:etr-reduction-spec}: the architecture,
activation function, loss/threshold convention, permitted trainable
parameters, and exact acceptance predicate must match, or a polynomial
model-conversion reduction must be given.

\subsection{Required gadget ledger}

The reduction must provide the following constructions with forward and
reverse correctness proofs.

\begin{center}
\begin{tabular}{p{0.22\linewidth}p{0.50\linewidth}p{0.14\linewidth}}
\toprule
Object & Required invariant & Status\\
\midrule
Rational constant & A designated signal/parameter equals the encoded rational. & Open\\
Signed variable & An unrestricted real variable is represented despite nonnegative ReLU outputs. & Open\\
Equality & Two represented expressions are equal iff the exact output constraints hold. & Open\\
Addition & A represented output equals the sum of represented inputs. & Open\\
Multiplication & A represented output equals the product of two represented unknowns, with legal sharing in the network model. & Open\\
Weak/strict sign & The source sign predicate is preserved using only permitted exact constraints. & Open\\
Conjunction & All source constraints hold simultaneously under one shared assignment. & Open\\
Reuse & Multiple occurrences of one source variable denote one real value without exponential copying. & Open\\
\bottomrule
\end{tabular}
\end{center}

Unknown weights multiply activations in an affine layer, so the network
equations contain bilinear terms.  Turning that observation into a multiplication
gadget requires more than noting the term's existence: both operands must be
represented as source variables, their signs must be handled, and all uses
must remain synchronized by the permitted parameter-sharing rules.

\subsection{Size and bit bounds}

For a valid many-one reduction, the construction must satisfy
\begin{equation}
  |\langle I_\varphi\rangle|
  \le \poly(|\langle\varphi\rangle|).
\end{equation}
This includes the number of nodes, samples, trainable incidences, local state
variables, and the bit lengths of rational constants.  Correctness must be
proved in both directions; a gadget that introduces unintended ReLU branches
is insufficient.

\begin{problem}[Model-matching ETR reduction]
\label{prob:etr-reduction}
Prove \cref{eq:etr-reduction-spec} for the exact target-equality language of
\cref{sec:problem}, with polynomial construction size and coefficient bit
length, or modify the language and every upstream theorem in a stated,
consistent manner.
\end{problem}

\begin{remark}[Current status]
Membership in \(\ExistsR\) is proved in \cref{prop:nntp-in-existsr}.  Hardness
for this exact formal language remains \openstatus until
\cref{prob:etr-reduction} is closed.
\end{remark}
